\section{Introduction}
\label{sec:introduction}

Remote-sensing change detection (CD) estimates a binary change mask from two co-registered images acquired over the same area at different times. The task underpins urban expansion monitoring, disaster assessment, land-cover analysis, and infrastructure management. In practice, binary CD is difficult because the model must compare two observations while being invariant to nuisance variation from illumination shifts, seasonal changes, shadows, sensor differences, and small registration errors. These factors are particularly visible near object boundaries, where a model can correctly identify the approximate changed region and still produce a poor mask if the predicted boundary is shifted, fragmented, or over-smoothed.

Deep CD methods have evolved along three main directions. Convolutional Siamese networks compare bi-temporal images through early fusion or feature differencing, as in FC-EF, FC-Siam-conc, and FC-Siam-diff~\cite{daudt2018fully}. Later CNN encoder--decoder models improve localization through skip connections, dense feature reuse, and multi-scale attention~\cite{ronneberger2015unet,fang2022snunet,chen2021dasnet}. These models are effective because convolutions provide useful locality and translation bias, but their context modeling is bounded by the receptive field, which can limit reasoning about spatially distributed changes. Transformer-based methods address long-range interaction more directly: BIT models bi-temporal tokens with a Transformer encoder~\cite{chen2022bit}, and ChangeFormer uses a hierarchical Siamese Transformer for dense CD~\cite{bandara2022changeformer}. Wijenayake et al.~\cite{precision_spatiotemporal_fusion_iciis2025} further demonstrate that precise spatio-temporal feature fusion yields gains over standard differencing strategies. Despite the progress, attention-based models carry high memory and compute cost at dense resolution, and stronger global representations do not automatically resolve boundary localization errors.

Recent visual state-space models (SSMs) provide a different route. Mamba introduces selective state-space sequence modeling with linear complexity~\cite{gu2023mamba}, and VMamba adapts it to vision through two-dimensional selective scanning~\cite{liu2024vmamba}. MambaVision develops a hybrid Mamba--Transformer backbone for visual recognition and dense prediction tasks~\cite{hatamizadeh2025mambavision}. Mamba-based CD models, including ChangeMamba~\cite{chen2024changemamba}, HAM-CD~\cite{li2026hamcd}, and Mamba-CD~\cite{peng2026mambacd}, show that SSM backbones are competitive for binary and semantic change detection. Wijenayake et al.\ in Mamba-FCS~\cite{mambafcs} further demonstrate that combining frequency-domain fusion with change-guided attention and a SeK-inspired loss improves semantic CD. However, applying a Mamba-based encoder alone is not a complete solution. A CD model must still construct temporal evidence from both images and preserve fine boundaries in the output mask. A controlled benchmark of SSM backbones for remote-sensing segmentation by Wasalathilaka et al.~\cite{wasalathilaka2026controlled} found that boundary delineation remains the dominant failure mode under domain shift, suggesting that encoder scaling alone is insufficient for boundary-sensitive tasks.

We propose \textbf{\ours}, a MambaVision-based architecture for binary remote-sensing CD. The model builds temporal features using raw paired features, absolute difference, and signed difference, giving the decoder access to both magnitude-based and direction-aware evidence. The core novelty is \textbf{\drbi}, which applies two separate learned gating mechanisms to temporal features before decoding: a region gate that produces stable, coherent changed-area responses, and a Sobel-conditioned boundary gate that selects contour-sensitive evidence. These two pathways are combined through an adaptive receptive-field decoder and refined by a bounded residual head that applies a controlled correction to coarse logits near boundaries.

The contributions of this paper are:
\begin{itemize}
    \item We present \textbf{\ours}, a MambaVision-based binary CD architecture integrating shared bi-temporal encoding, explicit four-stream temporal feature construction, and boundary-aware refinement.
    \item We propose \textbf{\drbi}, a lightweight module that decouples temporal evidence into complementary region-aware and boundary-aware gated pathways before multi-scale decoding.
    \item We integrate adaptive receptive-field decoding, CRAM-lite spatial modulation, and bounded residual boundary refinement into a unified pipeline, demonstrating that each component contributes measurably.
    \item We report verified results on DSIFN-CD and WHU-CD, averaged over independent runs, and present a controlled ablation study over seven variants (A0--A6) that isolates each component's effect on F1, IoU, and OA.
\end{itemize}
